\section{Discussion and limitations}
\label{sec:discussion}

\paragraph{What the scores mean.}
Every score here is a distance to a regularized solution of an underdetermined problem. No ground-truth DEM exists for these images, so neither our results nor the baseline's show that a DEM is physically right. A label-free model inherits the solver's regularizer through its objective. It emulates the solver's choice of answer and makes no independent claim about the Sun.

\paragraph{Coverage.}
The test window has no flare far above X1, and we did not treat saturated pixels specially, so behavior on large X-class events is untested. Training reads one solver output, the flag that drops the 0.01\% to 0.02\% of pixels the solver could not solve at all. We trained the two 1.5M configurations with three seeds each and the others once. We treat AIA alone, without Hinode/XRT. Our output is a single DEM, whereas DEMNet also predicts a distribution per bin, which we do not attempt.

\paragraph{The remaining gap.}
Our remaining flare-core error is in the height of the tallest peaks. Above a solver peak of 3,200 units our model reaches 0.53 of the solver's peak and the baseline 0.70, and those 269 pixel-targets hold 83\% of our squared error. Part of this is training: at validation flare cores the final model's median amortization gap is 21\%, against 5\% to 12\% elsewhere, and it reproduces 87\% of 94\Ang{} where the minimizer reaches 93\%.

\paragraph{Speed.}
On one NVIDIA B200 GPU the 1.5M and 409k networks take 0.32 and 0.11 seconds per $2048\times2048$ image, or 13.3M and 36.7M pixels per second. On the 16 cores of a solver job they take 36 and 19 seconds. BP took 221 seconds of solving on those cores before any retries, which makes the GPU about 700 times faster for the 1.5M model and the same CPU cores 6 and 12 times faster. Both pipelines still pay about 4 minutes of CPU deconvolution per image, so the end-to-end gain over a 9.5-minute BP job stays near a factor of two until that step also moves to a GPU.

\paragraph{Directions.}
\Cref{prop:law} suggests two refinements, neither trained yet. A per-pixel, per-channel weight proportional to $\lambda_c h_c$, with the single-edge price standing in for $\lambda_c$, needs only $R$, $B$ and $\sigma$ and would hold pixels to a roughly common tolerance. An $\ell_1$ exterior penalty is exact~\cite{han1979exact}: once $\mu$ exceeds $\lambda_c h_c$ in every binding channel (about 24 at the flare core of \cref{sec:law}), its minimizer solves the LP itself, at the cost of a slope that jumps at the band edge.

\section{Conclusion}
\label{sec:conclusion}

Training on a solver's objective is enough to emulate it at full-disk scale. Our per-pixel models never see a solver DEM and hold their own against a supervised baseline scored on the same pixels, although that baseline keeps DEM MSE and AIA MAE on BP and our means remain sensitive to a few pixels. The difficult part was the loss. A hard noise band turned into a penalty behaves as a price, and the LP says what that price is: \cref{eq:law} gives how far outside its band each channel settles, and \cref{prop:law} gives the weight below which that stops being true. At still lower weights hot plasma is lost, which explains a failure that training could not fix. More generally, when a network is trained on an objective rather than on labels, that objective can be solved directly on a sample of inputs, which is a cheap way to learn whether a failure lies in the loss or in the training. We recommend doing this before scaling up.
